\subsection{Image Classification with Linear Classifiers}\label{sec:image-classification-with-linear-classifiers}

\subsubsection{The Image Classification Problem}\label{sec:the-image-classification-problem}

Until now, we've been discussing \textbf{how to represent images} and \textbf{how to process them using correlation filters}. Now the question becomes: \emph{\textbf{given an image, how can a computer automatically recognize what it contains?}}

\highspace
\textcolor{Green3}{\faIcon{bullseye} \textbf{Motivation.}} Humans solve image recognition effortlessly. We recognize objects almost instantly, regardless of: viewpoint, lighting, scale, background, partial occlusions. For a computer, however, an image is \textbf{just a collection of numbers}. The goal of \textbf{image classification} is to teach the computer to associate these numbers with a semantic label.

\highspace
Formally, suppose we have an RGB image:
\begin{equation*}
    I \in \mathbb{R}^{R \times C \times 3}
\end{equation*}
The classifier must assign it to \textbf{one category} among a predefined set $\Lambda$ of $K$ possible categories:
\begin{equation*}
    \Lambda = \{ \lambda_1, \lambda_2, \ldots, \lambda_K \}
\end{equation*}
For example:
\begin{equation*}
    \Lambda = \left\{\text{wheel}, \text{car}, \text{castle}, \text{baboon}, \dots\right\}
\end{equation*}
So the task is simply: image $\to$ classifier $\to$ label.

\highspace
\begin{flushleft}
    \textcolor{Red2}{\faIcon{exclamation-triangle} \textbf{Classification is more than choosing a label}}
\end{flushleft}
A modern classifier usually does \textbf{not} produce only one word. Instead, it computes a \textbf{confidence score} for every possible class. For example, a classifier might output:
\begin{equation*}
    \text{scores} = \left[0.1, 0.7, 0.2\right]
\end{equation*}
This means that the classifier is 10\% confident that the image is a wheel, 70\% confident that it is a car, and 20\% confident that it is a castle. The predicted class is simply the one with the \textbf{highest score}.

\highspace
Mathematically, we can define the \textbf{classification task} as:
\begin{equation*}
    \mathcal{K}: I \to y
\end{equation*}
Where $I$ is the \textbf{input image}, $y$ is the \textbf{predicted label}, and $\mathcal{K}$ is the \textbf{classifier}. Since $y \in \Lambda$, the \hl{classifier is simply a function that maps an image into one of the predefined classes.} Initially, the classifier takes random numbers $\mathcal{K}_{\theta}$ as parameters $\theta$. After training, these parameters are adjusted to \hl{encode everything the model has learned about recognizing different classes.}

\highspace
\textcolor{Green3}{\faIcon{tools} \textbf{Output Space.}} If there are $L$ possible classes, the classifier actually computes $L$ scores simultaneously. Thus, instead of producing a single label:
\begin{equation*}
    \mathcal{K}: I \to \Lambda
\end{equation*}
We can think of the classifier as producing a vector of scores:
\begin{equation*}
    \mathcal{K}: \mathbb{R}^{R \times C \times 3} \to \left[0,1\right]^{L}
\end{equation*}
Meaning that the input is an RGB image, and the \textbf{output is a vector} of $L$ \textbf{confidence values} (often probabilities after a Softmax layer). Compared to a single label, the \hl{scores reveal something important}: the \hl{classifier can express \textbf{uncertainty}.} For example, it might be 70\% sure that an object is a car and 30\% sure that it is a truck. This is crucial for real-world applications, in which decisions must be made based on confidence levels.

\highspace
\begin{flushleft}
    \textcolor{Red2}{\faIcon{question-circle} \textbf{What is the Image Classification Problem?}}
\end{flushleft}
The \definition{Image Classification Problem} can be stated as: \hl{given an image, automatically determine which semantic category it belongs to.} Formally, given $I \in \mathbb{R}^{R \times C \times 3}$, find $y \in \Lambda$ where $I$ is the image, $\Lambda$ is the finite set of possible classes, and $y$ is the correct class label.

\highspace
\textcolor{Red2}{\faIcon{question-circle} \textbf{Why is this actually difficult?}} If images were always identical, classification would be trivial. Unfortunately, the same object can appear in infinitely many ways. For example, the class ``\emph{Dog}'' may contain images like \emph{a dog running in a park}, \emph{a dog sitting on a couch}, \emph{a dog playing with a ball}, etc. All of these belong to the \textbf{same class}. So the classifier must ignore variations that \textbf{do not change the object's identity}.

\highspace
But this is \textbf{not the only challenge}. Different classes may look very similar. For example, a \emph{Wolf}, a \emph{Dog}, and a \emph{Fox} or a \emph{Car}, a \emph{Truck}, and a \emph{SUV} may share many visual features. The classifier must therefore learn:
\begin{itemize}
    \item \textbf{Intra-class invariance}: The ability to recognize the same object despite appearance changes (e.g., different angles, lighting, or occlusions).
    \item \textbf{Inter-class discrimination}: The ability to distinguish between visually similar objects (e.g., a dog vs. a wolf).
\end{itemize}

\highspace
\textcolor{Green3}{\faIcon{question-circle} \textbf{Why do we need Machine Learning?}} We could imagine writing rules manually (e.g., ``if the image has four wheels, it is a car''). However, this approach is impractical for many reasons:
\begin{itemize}
    \item The number of possible rules is enormous, and it is impossible to cover all cases.
    \item The rules would be brittle and fail to generalize to new images.
    \item The rules would be hard to maintain and update as new classes are added.
\end{itemize}
The number of rules quickly becomes impossible to manage. Instead, we let a model \textbf{learn} these rules from examples.

\highspace
In general, the complete image classification process introduced in this section can be summarized as follows:
\begin{enumerate}
    \item \textbf{Input:} An RGB image $I \in \mathbb{R}^{R \times C \times 3}$.
    \item \textbf{Classifier:} A function $\mathcal{K}$ that takes the image as input with parameters $\theta$ and produces a vector of scores.
    \item \textbf{Output:} A vector of scores $\text{scores} \in \left[0,1\right]^{L}$, where each score corresponds to a class in the predefined set $\Lambda$.
    \item \textbf{Decision:} The predicted label $y$ is the class with the highest score, i.e., $y = \argmax_{i} \text{scores}[i]$.
\end{enumerate}

\highspace
But \textbf{\emph{how does the classifier compute these scores from millions of pixel values?}} The next subsection answers this by introducing the simplest possible classifier: the \emph{linear classifier}. Rather than searching for objects explicitly, it learns a set of weights that directly transform the image pixels into class scores. This will be the foundation for understanding neural networks and, later, convolutional neural networks.